% ID: SEL-D0072-OB-E02
% ESTADO: cotejada. Segunda pasada visual y borrador compilado revisados.
% Fuente: Selectividad/sel_2009_matematicas_aplicadas/mat apli CC SS06.pdf
% SHA-256: b371bdd6ba8b176301371f25945c7fb030327ff4501a6a02aa6f109a3a4f53e1
% Página: 2; opción B; ejercicio original 2.

Sea la función $f:\mathbb{R}\to\mathbb{R}$ definida mediante
\[
f(x)=\begin{cases}e^{-x}&\text{si }x\leq0\\x^3-x+1&\text{si }x>0\end{cases}.
\]
\begin{enumerate}[label=\alph*)]
\item \textbf{(1 punto)} \textquestiondown{}Es $f$ continua en $x=0$? \textquestiondown{}Es continua en su dominio?
\item \textbf{(1 punto)} \textquestiondown{}Es $f$ derivable en $x=0$? \textquestiondown{}Es derivable en su dominio?
\item \textbf{(1 punto)} Halle la ecuación de la recta tangente a la gráfica de $f$ en el punto de abscisa $x=1$.
\end{enumerate}
